\chapter{Bit Manipulation}

\section{Binary number}
A binary number is a number expressed in the base-2 numeral system or binary numeral system, it is a method of mathematical expression which uses only two symbols: typically ``0" (zero) and ``1" (one).

We say that a certain bit is \textbf{set}, if it is one, and cleared if it is zero.

\section{Bitwise Operator}

\vspace{3ex}
\textbf{\large{Bitwise AND}: }
Returns 1 only if both the bits are 1.

\begin{tabularx}{9cm}{|P{1cm}|C|C|C|C|C|C|C|C|}\hline
   a = 25 &
   0 & 0 & 0 & 1 & 1 & 0 & 0 & 1\\\hline
   b = 14 &
   0 & 0 & 0 & 0 & 1 & 1 & 1 & 0\\\hline
   a \& b &
   0 & 0 & 0 & 0 & 1 & 0 & 0 & 0\\\hline
\end{tabularx}

\vspace{3ex}
\textbf{\large{Bitwise OR}: }
Returns 1 if one of the bits is 1.

\begin{tabularx}{9cm}{|P{1.5cm}|C|C|C|C|C|C|C|C|}\hline
   a = 25 &
   0 & 0 & 0 & 1 & 1 & 0 & 0 & 1\\\hline
   b = 14 &
   0 & 0 & 0 & 0 & 1 & 1 & 1 & 0\\\hline
   a | b &
   0 & 0 & 0 & 1 & 1 & 1 & 1 & 1\\\hline
\end{tabularx}

\vspace{3ex}
\textbf{\large{Bitwise NOT}: }
Returns 1 if one of the bits is 1.

\begin{tabularx}{9cm}{|P{1.5cm}|C|C|C|C|C|C|C|C|}\hline
   a = 25 &
   0 & 0 & 0 & 1 & 1 & 0 & 0 & 1\\\hline
   $\sim$ a &
   1 & 1 & 1 & 0 & 0 & 1 & 1 & 0\\\hline
\end{tabularx}

\vspace{3ex}
\textbf{\large{Bitwise XOR}: }
Returns 0 if both the bits are same else 1

\begin{tabularx}{9cm}{|P{1.5cm}|C|C|C|C|C|C|C|C|}\hline
   a = 25 &
   0 & 0 & 0 & 1 & 1 & 0 & 0 & 1\\\hline
   b = 14 &
   0 & 0 & 0 & 0 & 1 & 1 & 1 & 0\\\hline
   a $\wedge$ b &
   0 & 0 & 0 & 1 & 0 & 1 & 1 & 1\\\hline
\end{tabularx}

\vspace{3ex}
\textbf{\large{Left Shift}: }
Shifts a towards left by n digits.
\texttt{a << b = a * $\texttt{2}^\texttt{b}$}

\begin{tabularx}{9cm}{|P{1.5cm}|C|C|C|C|C|C|C|C|}\hline
   \rowcolor{string!20}
   \cellcolor{white}
   a = 25 &
   0 & 0 & 0 & 1 & 1 & 0 & 0 & 1\\\hline
   \rowcolor{gray!20}
   \cellcolor{white}
   a \texttt{<<} 2 &
   0 & 1 & 1 & 0 & 0 & 1 & 
   \cellcolor{white}0 & \cellcolor{white}0\\\hline
\end{tabularx}

\vspace{3ex}
\textbf{\large{Right Shift}: }
Shifts a towards right by n digits.
\texttt{a >> b = $\frac{\texttt{a}}{\texttt{2}^\texttt{b}}$}

\begin{tabularx}{9cm}{|P{1.5cm}|C|C|C|C|C|C|C|C|}\hline
   \rowcolor{string!20}
   \cellcolor{white}
   a = 25 &
   0 & 0 & 0 & 1 & 1 & 0 & 0 & 1\\\hline
   \rowcolor{gray!20}
   \cellcolor{white}
   a \texttt{>>} 2 &
   \cellcolor{white}0 & \cellcolor{white}0 &
   0 & 0 & 0 & 1 & 1 & 0\\\hline
\end{tabularx}

\vspace{2ex}
\section{Bit Masking}


\begin{tabular}{lll}
   x \ \& \ 1 & : &
   Evaluates to 1 if the number is odd else evaluates to 0.
   \\
   x \ \& \ (x-1) & : & Clears the rightmost set bit of x.
   \\
   x \ \& \ $\sim$(x-1) & : &
   Isolates the rightmost set bit of x.
   \\
   \\
   1 \ \texttt{<<} \ i & : & Represent a number with only the $i$-th bit set.
   \\
   $\sim$(1 \ \texttt{<<} \ i) & : & Represent a number with all bits set except the $i$-th bit.
   \\
   \\
   n \& (1 \texttt{<<} i) & : & Checks whether $i$-th bit  in $n$ is set or not.
   \\
   n \& $\sim$(1 \texttt{<<} i) & : & Clears the $i$-th bit in $n$.
   \\
   n | (1 \texttt{<<} i) & : & Sets the $i$-th bit  in $n$ without affecting other bits.
   \\
   n | $\sim$(1 \texttt{<<} i) & : & Clears all the bits except $i$-th bit in $n$.
   \\
   n $\wedge$ (1 \texttt{<<} i) & : & Flips the $i$-th bit in $n$ without affecting other bits.
   \\
   n $\wedge$ $\sim$(1 \texttt{<<} i) & : & Flips the $i$-th bit and complements all other bits in $n$.
   \\
\end{tabular}

\pagebreak
\subsection{Get, Set, Clear \codeit{i}-th bit}
\textbf{Set \codeit{i}-th bit}
\begin{code}[language=C++]
   int set_ith_bit(int n, int i){
      int mask = (1 << i);
      return (n|mask);
   }
\end{code}


\textbf{Clear \codeit{i}-th bit}
\begin{code}[language=C++]
   int clear_ith_bit(int n, int i){
      int mask = ~(1 << i);
      return (n & mask);
   }
\end{code}


\textbf{Update \codeit{i}-th bit}
\begin{code}[language=C++]
   int update_ith_bit(int n, int i, int v){
      clearIthBit(n,i);
      int mask = (v << i);
      return (n|mask);
   }
\end{code}


\textbf{Clear last \codeit{i} bits}
\begin{code}[language=C++]
   int clear_last_i_bit(int n, int i){
      int mask = (-1 << i);
      return (n & mask);
   }
\end{code}


\textbf{Clear bits in range}
\begin{code}[language=C++]
   int clear_bits_in_range(int n, int i, int j){
      int a = -1 << (j+1),  b = (1 << i) - 1;
      int mask = (a|b);
      return (n & mask);
   }
\end{code}


\textbf{Update bits in range}
\begin{code}[language=C++]
   int update_bits_in_range(int n, int i, int v){
      clear_bits_in_range(n,i);
      int mask = (v << i);
      return (n|mask);
   }
\end{code}


\textbf{Check if an integer is a power of 2}
\begin{code}[language=C++]
   bool isPowerOfTwo(unsigned int n){
      return !(n & (n-1));
   }
\end{code}


\textbf{Count set bits} \tab O(number of set bits)
\begin{code}[language=C++]
   int countSetBits(int n){
      int count = 0;
      while (n){
         n = n & (n - 1);
         count++;
      }
      return count;
   }
\end{code}


\textbf{Find the unique number from an array of 2N + 1 numbers.}\\
\begin{tabular}{ll}      
   Principle &:\tab  \texttt{a $\wedge$ a = 0}
   \tabs  \texttt{0 $\wedge$ a = a}\\
   Approach &: Take Commulative XOR\\
   Time complexity &: O(N)
\end{tabular}